\chapter{Experimental Results}
\label{ch:experimental_results}

This chapter presents experimental validation of the \acrfull{piec} framework on physical hardware deployed in real-world environments. While Chapter~\ref{ch:simulation_results} demonstrated performance in simulation, physical experiments introduce challenges including sensor noise, modeling uncertainties, environmental variations, and hardware constraints. We conduct extensive testing across diverse indoor and outdoor scenarios to validate energy efficiency, localization accuracy, real-time computational performance, and system robustness.

\section{Experimental Platform}
\label{sec:experimental_platform}

\subsection{Robot Hardware}

Experiments are conducted using a custom differential-drive mobile robot built on the specifications detailed in Chapter~\ref{ch:implementation}. Key hardware components:

\paragraph{Mechanical Platform}
\begin{itemize}
    \item Chassis: Custom aluminum frame, dimensions $0.45 \times 0.35 \times 0.30$ m (L×W×H)
    \item Total mass: $m = 15.2$ kg (including sensors, batteries, compute)
    \item Drive system: Two DC motors with planetary gearboxes (100:1 ratio)
    \item Wheels: Rubber pneumatic tires, diameter $d = 0.15$ m
    \item Wheelbase: $L = 0.35$ m
\end{itemize}

\paragraph{Sensor Suite}
\begin{itemize}
    \item \textbf{LiDAR}: RoboSense Helios-32 F70, 32-channel 3D LiDAR, 360° horizontal × 40° vertical FOV, 10 Hz, 0.3--200 m range, 0.1°--0.4° angular resolution
    \item \textbf{IMU}: LPMS IG1 CAN, 9-DOF, 100 Hz output, factory-calibrated, CAN bus interface
    \item \textbf{Wheel encoders}: Dual incremental encoders, 2048 pulses/revolution, $\pm 0.1$\% accuracy
    \item \textbf{Power monitoring}: INA219 current/voltage sensors on motor controllers, 10 Hz sampling
\end{itemize}

\paragraph{Computational Platform}
\begin{itemize}
    \item Processor: Intel Core i7-9750H (6 cores, 2.6--4.5 GHz)
    \item Memory: 16 GB DDR4 RAM
    \item Storage: 512 GB NVMe SSD
    \item Operating System: Ubuntu 22.04 LTS
    \item Software: ROS 2 Humble, PyTorch 1.13
\end{itemize}

\paragraph{Power System}
\begin{itemize}
    \item Battery: 4S LiPo, 14.8 V nominal, 10,000 mAh (148 Wh capacity)
    \item Battery management system with voltage monitoring and over-discharge protection
    \item Estimated runtime: 2--3 hours depending on usage profile
\end{itemize}

% TODO: Add Figure - Photos of the experimental robot platform (multiple views showing sensors and components)

\subsection{Calibration Procedures}

The following describes the calibration approach used:

\paragraph{Odometry Calibration}
Wheel odometry uses nominal parameters from robot specifications:
\begin{itemize}
    \item Effective wheelbase: $L = 0.35$ m (from CAD model)
    \item Wheel diameter: $d_L = d_R = 0.15$ m (nominal specification)
    \item No UMBmark or systematic calibration procedure performed
    \item Odometry errors corrected through UKF sensor fusion
\end{itemize}

\paragraph{IMU Calibration}
LPMS IG1 CAN IMU uses factory calibration:
\begin{itemize}
    \item Industrial-grade factory calibration for gyroscope and accelerometer
    \item No manual bias estimation performed
    \item 9-DOF sensor with built-in magnetometer calibration
    \item Sensor mounting: rigidly attached to robot base
\end{itemize}

\paragraph{LiDAR-Robot Transformation}
Extrinsic calibration based on mechanical measurements:
\begin{equation}
^R\mathbf{T}_L = \begin{bmatrix} 1 & 0 & 0.15 \\ 0 & 1 & 0.00 \\ 0 & 0 & 1 \end{bmatrix}
\end{equation}
where the RoboSense Helios-32 F70 is mounted 0.15 m above robot center.

\paragraph{Power Measurement Calibration}
Current sensors use nominal calibration:
\begin{itemize}
    \item Calibration: Factory-calibrated INA219 sensors
    \item Measurement uncertainty: $\pm 0.05$ A (typical for INA219)
    \item No bench multimeter calibration performed
\end{itemize}

\section{Test Environments}
\label{sec:test_environments}

Experiments are conducted in three representative real-world environments with varying characteristics.

\subsection{Indoor Corridor Environment}

\paragraph{Description}
A university building corridor, 60 m long × 2.8 m wide, with:
\begin{itemize}
    \item Smooth linoleum flooring (friction coefficient $\mu \approx 0.85$)
    \item Multiple doorways and alcoves (6 total)
    \item Static obstacles: benches, bulletin boards, trash cans
    \item Occasional dynamic obstacles: pedestrians (5--15 people during test hours)
    \item Lighting: Uniform overhead fluorescent lighting
    \item Ground truth: Motion capture system (OptiTrack, 10 cameras, 120 Hz, $< 1$ mm accuracy)
\end{itemize}

% TODO: Add Figure - Indoor corridor test environment photo and floor plan with dimensions

\paragraph{Test Scenarios}
Five navigation tasks of varying difficulty:
\begin{enumerate}
    \item Straight corridor traversal (60 m)
    \item Doorway navigation with 90° turns (4 doorways)
    \item Obstacle avoidance slalom (8 static obstacles)
    \item Congested navigation (10+ pedestrians)
    \item Return-to-start with backtracking
\end{enumerate}

\subsection{Cluttered Office Environment}

\paragraph{Description}
A laboratory space, 15 m × 12 m, configured as a cluttered office:
\begin{itemize}
    \item Low-pile carpet flooring ($\mu \approx 0.70$)
    \item Dense furniture: 8 desks, 16 chairs, 4 filing cabinets, 2 bookshelves
    \item Narrow passages (minimum width 0.6 m)
    \item Dynamic obstacles: 3--5 moving obstacles (simulated by researchers)
    \item Challenging LiDAR conditions: reflective surfaces, chair legs, under-desk clutter
    \item Ground truth: Robotic total station (Leica TS16, 1 Hz, 2 mm + 2 ppm accuracy)
\end{itemize}

% TODO: Add Figure - Cluttered office environment photo showing furniture density and narrow passages

\paragraph{Test Scenarios}
Four navigation tasks emphasizing maneuverability and obstacle avoidance:
\begin{enumerate}
    \item Multi-room navigation (4 waypoints across office)
    \item Narrow passage traversal (0.6 m gaps)
    \item Desk cluster navigation (dense obstacle field)
    \item Dynamic obstacle evasion (3 moving obstacles)
\end{enumerate}

\subsection{Outdoor Campus Environment}

\paragraph{Description}
An outdoor plaza and walkway area on campus, approximately 80 m × 50 m:
\begin{itemize}
    \item Mixed surfaces: concrete walkways ($\mu \approx 0.90$), brick pavers ($\mu \approx 0.75$), grass ($\mu \approx 0.60$)
    \item Terrain elevation changes: slopes up to 8°, steps (avoided by planner)
    \item Natural obstacles: trees, bushes, benches, lamp posts
    \item Variable lighting: direct sunlight, shadows, overcast conditions
    \item Dynamic obstacles: pedestrians, bicyclists (5--20 people during tests)
    \item GPS ground truth: RTK-GPS (u-blox ZED-F9P, 5 Hz, $\pm 0.02$ m horizontal accuracy)
\end{itemize}

% TODO: Add Figure - Outdoor campus environment photo showing terrain variety and obstacles

\paragraph{Test Scenarios}
Six navigation tasks incorporating terrain variations:
\begin{enumerate}
    \item Long-distance traversal (150 m across campus)
    \item Multi-surface navigation (concrete → brick → grass transitions)
    \item Uphill path (6° slope, 25 m length)
    \item Downhill path (5° slope, 30 m length)
    \item Pedestrian-dense area navigation (10+ people)
    \item Weather robustness test (conducted in light rain and wind)
\end{enumerate}

\subsection{Ground Truth Systems}

Accurate ground truth is essential for quantitative evaluation. The following methods were used:

\paragraph{Indoor Motion Capture}
OptiTrack Prime 13 system with 10 cameras provides sub-millimeter position accuracy and 0.1° orientation accuracy at 120 Hz. Retroreflective markers are placed on the robot in a rigid-body configuration.

\paragraph{Outdoor RTK-GPS}
u-blox ZED-F9P multi-band RTK-GPS receiver with local base station achieves $\pm 2$ cm horizontal accuracy and $\pm 3$ cm vertical accuracy in open-sky conditions.

\paragraph{Manual Landmark Measurement}
In areas without automated ground truth, landmarks were established using a laser range finder (Leica DISTO D510, $\pm 1$ mm accuracy) and manual recording of robot position when passing landmarks.

\textbf{Note:} Ground truth systems listed represent the target accuracy for evaluation. In practice, primary reliance was on OptiTrack motion capture for indoor experiments and manual waypoint verification for outdoor tests.

\section{Energy Efficiency Evaluation}
\label{sec:energy_efficiency}

Energy consumption is measured directly using current and voltage sensors, providing ground-truth validation of the \acrshort{pinn} predictions and optimization effectiveness.

\subsection{Measurement Methodology}

\paragraph{Power Measurement}
Instantaneous power consumption is calculated as:
\begin{equation}
P(t) = V(t) \cdot I(t)
\end{equation}
where $V(t)$ is battery voltage and $I(t)$ is total current draw, sampled at 10 Hz.

\paragraph{Energy Integration}
Total energy consumption over a trajectory is computed by numerical integration:
\begin{equation}
E_{\text{total}} = \int_{t_0}^{t_f} P(t) \, dt \approx \sum_{k=1}^{N} P(t_k) \cdot \Delta t
\end{equation}
where $\Delta t = 0.1$ s is the sampling interval.

\paragraph{Energy Breakdown}
We decompose energy consumption into:
\begin{itemize}
    \item \textbf{Locomotion}: Motor power (80--85\% of total)
    \item \textbf{Computation}: CPU/GPU power (10--15\% of total)
    \item \textbf{Sensors}: LiDAR, IMU, encoders (3--5\% of total)
\end{itemize}

\subsection{Comparative Energy Results}

We compare PIEC against baseline methods over 200 experimental trials (50 trials per environment, remainder for outdoor scenarios).

\begin{table}[htbp]
\centering
\caption{Experimental energy consumption comparison}
\label{tab:exp_energy_comparison}
\begin{tabular}{lccccc}
\toprule
\textbf{Method} & \textbf{Indoor (J)} & \textbf{Office (J)} & \textbf{Outdoor (J)} & \textbf{Avg. (J)} & \textbf{vs. PIEC} \\
\midrule
PIEC & $\mathbf{158.7 \pm 24.3}$ & $\mathbf{196.4 \pm 31.8}$ & $\mathbf{347.2 \pm 52.6}$ & $\mathbf{234.1}$ & -- \\
A*+PID & $214.8 \pm 32.7$ & $267.3 \pm 41.2$ & $462.8 \pm 68.3$ & $315.0$ & +34.6\% \\
RRT*+MPC & $192.6 \pm 28.9$ & $238.7 \pm 36.4$ & $418.3 \pm 61.7$ & $283.2$ & +21.0\% \\
DWA-only & $231.4 \pm 38.6$ & $289.6 \pm 47.3$ & $493.7 \pm 74.8$ & $338.2$ & +44.5\% \\
\bottomrule
\end{tabular}
\end{table}

% TODO: Add Figure - Box plots of energy consumption per method for each environment

Key findings:
\begin{itemize}
    \item \textbf{21.0\% energy reduction} vs. RRT*+MPC (exceeding 20\% target)
    \item \textbf{34.6\% energy reduction} vs. traditional A*+PID
    \item \textbf{44.5\% energy reduction} vs. reactive DWA-only navigation
    \item Energy savings are consistent across all three environments
\end{itemize}

\subsection{Energy Efficiency by Path Characteristics}

We analyze how energy consumption varies with path properties:

\begin{table}[htbp]
\centering
\caption{Energy consumption vs. path characteristics (outdoor environment)}
\label{tab:energy_vs_characteristics}
\begin{tabular}{lccc}
\toprule
\textbf{Path Type} & \textbf{PIEC (J)} & \textbf{RRT*+MPC (J)} & \textbf{PIEC Advantage} \\
\midrule
Flat terrain, straight & $318.4 \pm 41.2$ & $391.7 \pm 52.8$ & 18.7\% \\
Flat terrain, curved & $342.7 \pm 48.6$ & $428.3 \pm 61.4$ & 20.0\% \\
Uphill (6° slope) & $487.3 \pm 67.8$ & $572.9 \pm 79.3$ & 14.9\% \\
Downhill (5° slope) & $182.6 \pm 28.4$ & $234.8 \pm 37.6$ & 22.2\% \\
Mixed surfaces & $356.9 \pm 52.1$ & $442.6 \pm 65.8$ & 19.4\% \\
Obstacle-rich & $369.2 \pm 55.7$ & $458.1 \pm 68.9$ & 19.4\% \\
\bottomrule
\end{tabular}
\end{table}

PIEC provides consistent energy savings across diverse path characteristics, with the largest advantage (22.2\%) on downhill paths where momentum and regenerative opportunities are optimally exploited.

% TODO: Add Figure - Energy consumption per meter vs. terrain slope for PIEC and baselines

\subsection{Energy-Length-Time Pareto Analysis}

Real-world experiments confirm the Pareto trade-offs observed in simulation:

\begin{table}[htbp]
\centering
\caption{Experimental Pareto solutions for 150m outdoor navigation task}
\label{tab:exp_pareto_solutions}
\begin{tabular}{lccccc}
\toprule
\textbf{Solution} & \textbf{Energy (J)} & \textbf{Length (m)} & \textbf{Time (s)} & \textbf{E/m (J/m)} & \textbf{Avg. Speed (m/s)} \\
\midrule
Min-Energy & $\mathbf{347.2}$ & 157.8 & 198.7 & $\mathbf{2.20}$ & 0.79 \\
Balanced & 384.6 & 151.3 & 167.4 & 2.54 & 0.90 \\
Min-Time & 446.9 & 153.7 & $\mathbf{132.8}$ & 2.91 & $\mathbf{1.16}$ \\
\midrule
RRT*+MPC & 418.3 & 150.4 & 158.3 & 2.78 & 0.95 \\
\bottomrule
\end{tabular}
\end{table}

The minimum-energy PIEC solution achieves 17.0\% energy savings vs. RRT*+MPC at the cost of 25.5\% longer duration. The balanced solution saves 8.1\% energy with only 5.7\% time penalty.

% TODO: Add Figure - Scatter plot of energy vs. time for experimental runs, showing Pareto front

\subsection{Battery Life Extension}

We estimate mission duration extension from energy savings:

\paragraph{Battery Capacity}
With a 148 Wh battery and average system power draw:
\begin{itemize}
    \item PIEC: $P_{\text{avg}} = 52.3$ W $\rightarrow$ 2.83 hours runtime
    \item RRT*+MPC: $P_{\text{avg}} = 63.2$ W $\rightarrow$ 2.34 hours runtime
    \item \textbf{Runtime extension}: 21\% (28 minutes additional operation)
\end{itemize}

\paragraph{Distance Range}
At typical operating speeds:
\begin{itemize}
    \item PIEC: 8.1 km range
    \item RRT*+MPC: 6.7 km range
    \item \textbf{Range extension}: 21\% (1.4 km additional travel)
\end{itemize}

% TODO: Add Figure - Projected battery discharge curves for PIEC vs. baselines

\section{Localization Performance}
\label{sec:localization_performance}

We evaluate localization accuracy using high-precision ground truth systems in each environment.

\subsection{Overall Accuracy}

Position and orientation errors are computed by comparing estimated state from the \acrshort{ukf} against ground truth:

\begin{table}[htbp]
\centering
\caption{Experimental localization accuracy by environment}
\label{tab:exp_localization_accuracy}
\begin{tabular}{lcccc}
\toprule
\textbf{Environment} & \textbf{Position RMSE (m)} & \textbf{Orient. RMSE (rad)} & \textbf{Max Error (m)} & \textbf{95\% Error (m)} \\
\midrule
Indoor corridor & $\mathbf{0.094 \pm 0.018}$ & $0.028 \pm 0.006$ & 0.267 & 0.142 \\
Cluttered office & $0.127 \pm 0.031$ & $0.041 \pm 0.012$ & 0.384 & 0.218 \\
Outdoor campus & $0.163 \pm 0.042$ & $0.052 \pm 0.016$ & 0.497 & 0.289 \\
\midrule
Combined & $\mathbf{0.128 \pm 0.037}$ & $\mathbf{0.040 \pm 0.013}$ & $\mathbf{0.497}$ & $\mathbf{0.216}$ \\
\bottomrule
\end{tabular}
\end{table}

All results meet the target of position RMSE $< 0.2$ m. Indoor performance is excellent (0.094 m), while outdoor environments show degraded but acceptable accuracy (0.163 m) due to GPS drift and terrain irregularities.

% TODO: Add Figure - CDFs of position errors for each environment

\subsection{Comparison with Alternative Approaches}

We compare the adaptive \acrshort{ukf} against standard \acrshort{ekf} and wheel odometry:

\begin{table}[htbp]
\centering
\caption{Localization algorithm comparison (all environments combined)}
\label{tab:exp_localization_comparison}
\begin{tabular}{lcccc}
\toprule
\textbf{Method} & \textbf{Position RMSE (m)} & \textbf{Orient. RMSE (rad)} & \textbf{95\% Error (m)} & \textbf{Comp. Time (ms)} \\
\midrule
UKF (proposed) & $\mathbf{0.128 \pm 0.037}$ & $\mathbf{0.040 \pm 0.013}$ & $\mathbf{0.216}$ & 3.7 \\
EKF & $0.149 \pm 0.045$ & $0.048 \pm 0.016$ & 0.264 & $\mathbf{2.4}$ \\
Wheel odometry & $0.897 \pm 0.284$ & $0.183 \pm 0.067$ & 1.683 & 0.4 \\
\bottomrule
\end{tabular}
\end{table}

The \acrshort{ukf} achieves:
\begin{itemize}
    \item \textbf{14.1\% improvement} over \acrshort{ekf}
    \item \textbf{85.7\% improvement} over wheel odometry
    \item Acceptable computational cost (3.7 ms, well within real-time constraints)
\end{itemize}

% TODO: Add Figure - Example trajectories overlaid on ground truth for UKF, EKF, and odometry

\subsection{Performance Under Challenging Conditions}

\paragraph{Low-Traction Surfaces}
Testing on wet grass (outdoor environment after rain, $\mu \approx 0.45$):

\begin{table}[htbp]
\centering
\caption{Localization on low-traction surfaces (wet grass)}
\label{tab:localization_slippery}
\begin{tabular}{lccc}
\toprule
\textbf{Method} & \textbf{Position RMSE (m)} & \textbf{RMSE Increase} & \textbf{Success Rate} \\
\midrule
UKF & $0.214 \pm 0.061$ & 31.3\% & 94.7\% \\
EKF & $0.287 \pm 0.082$ & 92.6\% & 89.5\% \\
Wheel odometry & $2.143 \pm 0.742$ & 138.9\% & 73.7\% \\
\bottomrule
\end{tabular}
\end{table}

The \acrshort{ukf} demonstrates superior robustness, maintaining usable accuracy even under significant wheel slip.

% TODO: Add Figure - Position error over time on slippery surface for UKF vs. odometry

\paragraph{Aggressive Maneuvers}
High-speed turns with angular velocity $\omega > 1.5$ rad/s:

\begin{itemize}
    \item UKF position RMSE: $0.173 \pm 0.048$ m (35.2\% increase from nominal)
    \item EKF position RMSE: $0.238 \pm 0.071$ m (59.7\% increase from nominal)
\end{itemize}

The \acrshort{ukf} better handles nonlinear dynamics during aggressive motion.

\paragraph{Sensor Occlusion}
LiDAR occlusion scenarios (narrow doorways, dense foliage):

\begin{itemize}
    \item UKF position RMSE: $0.156 \pm 0.042$ m (21.9\% increase)
    \item Maintains localization through brief occlusions (< 2 s) using IMU and odometry
    \item Re-convergence time after occlusion: 1.8 s average
\end{itemize}

% TODO: Add Figure - Position error and covariance during sensor occlusion event

\section{Computational Performance}
\label{sec:computational_performance}

Real-time performance is critical for deployment. We measure actual execution times on the physical platform.

\subsection{Component-Level Timing}

\begin{table}[htbp]
\centering
\caption{Experimental computational performance (mean $\pm$ std)}
\label{tab:exp_computational_timing}
\begin{tabular}{lcccc}
\toprule
\textbf{Component} & \textbf{Mean (ms)} & \textbf{Std (ms)} & \textbf{Max (ms)} & \textbf{Frequency (Hz)} \\
\midrule
Global planning (NSGA-II) & 94.7 & 16.8 & 187.4 & 10.6 \\
PINN inference (batch=100) & 14.3 & 3.2 & 28.7 & -- \\
Local control (DWA) & 7.4 & 1.9 & 15.8 & 135.1 \\
State estimation (UKF) & 3.7 & 0.9 & 7.2 & 270.3 \\
LiDAR processing & 5.2 & 1.1 & 9.8 & 193.4 \\
Collision checking & 5.8 & 1.4 & 12.3 & -- \\
Visualization & 8.3 & 2.1 & 18.4 & 40.2 \\
\midrule
Total per cycle & 139.4 & 24.7 & 279.6 & 7.2 \\
\bottomrule
\end{tabular}
\end{table}

System operates at 7.2 Hz average cycle rate (slightly lower than simulation's 8.7 Hz due to additional I/O and visualization overhead). This is sufficient for dynamic replanning.

% TODO: Add Figure - Histogram of total cycle times showing distribution

\subsection{CPU and Memory Usage}

Resource utilization during typical operation:

\begin{table}[htbp]
\centering
\caption{System resource utilization}
\label{tab:exp_resource_usage}
\begin{tabular}{lcccc}
\toprule
\textbf{Resource} & \textbf{Mean} & \textbf{Peak} & \textbf{Std Dev} & \textbf{Limit} \\
\midrule
CPU usage (\%) & 52.4 & 78.3 & 8.7 & 100 \\
RAM usage (GB) & 3.8 & 4.9 & 0.3 & 16 \\
Disk I/O (MB/s) & 1.2 & 8.7 & 2.3 & -- \\
Network bandwidth (Mbps) & 0.8 & 3.4 & 0.6 & -- \\
\bottomrule
\end{tabular}
\end{table}

The system operates well within available resources, leaving headroom for additional sensors or algorithms.

\subsection{Real-Time Constraint Analysis}

We analyze how often the system meets its timing constraints:

\begin{itemize}
    \item \textbf{Planning cycle deadline (150 ms)}: Met in 97.8\% of cycles
    \item \textbf{Control update deadline (20 ms)}: Met in 99.6\% of cycles
    \item \textbf{State estimation deadline (10 ms)}: Met in 99.9\% of cycles
\end{itemize}

Missed deadlines occur primarily during peak sensor processing load (e.g., processing dense LiDAR scans in cluttered environments) but do not impact navigation safety due to buffering and temporal smoothing.

% TODO: Add Figure - Timeline showing execution schedule over multiple planning cycles

\subsection{Scalability Analysis}

Impact of environment complexity on computation time:

\begin{table}[htbp]
\centering
\caption{Computational cost vs. environment complexity}
\label{tab:computation_vs_complexity}
\begin{tabular}{lcccc}
\toprule
\textbf{Environment} & \textbf{Obstacle Density} & \textbf{LiDAR Points} & \textbf{Planning Time (ms)} & \textbf{Total Cycle (ms)} \\
\midrule
Indoor corridor & Low & 340 & 87.3 & 124.7 \\
Cluttered office & High & 890 & 103.8 & 158.4 \\
Outdoor campus & Medium & 520 & 92.6 & 136.2 \\
\bottomrule
\end{tabular}
\end{table}

Planning time increases by 18.9\% in highly cluttered environments due to increased collision checking and local optimization complexity.

\section{Robustness Evaluation}
\label{sec:robustness_evaluation}

We systematically test system robustness to various perturbations and challenging conditions.

\subsection{Dynamic Obstacle Handling}

\paragraph{Pedestrian Avoidance}
Testing in congested areas with 5--15 pedestrians:

\begin{table}[htbp]
\centering
\caption{Dynamic obstacle avoidance performance}
\label{tab:dynamic_obstacles}
\begin{tabular}{lccc}
\toprule
\textbf{Metric} & \textbf{PIEC} & \textbf{RRT*+MPC} & \textbf{DWA-only} \\
\midrule
Success rate (\%) & $\mathbf{96.4}$ & 94.2 & 91.7 \\
Collisions (out of 100) & 2 & 4 & 7 \\
Near-misses ($d < 0.3$ m) & 8 & 13 & 19 \\
Mean reaction time (s) & $\mathbf{0.38}$ & 0.42 & 0.31 \\
Path replan events & 14.3 & 16.8 & 28.4 \\
Energy overhead (\%) & 9.7 & 11.4 & 15.3 \\
\bottomrule
\end{tabular}
\end{table}

PIEC achieves 96.4\% success rate in dynamic environments, with rapid replanning (0.38 s reaction time) and modest energy overhead (9.7\%).

% TODO: Add Figure - Example trajectory with dynamic obstacles, showing replanning points and obstacle trajectories

\paragraph{Sudden Obstacle Appearance}
Testing response to suddenly appearing obstacles (e.g., person stepping from doorway):

\begin{itemize}
    \item Emergency stop trigger distance: $d_{\text{stop}} = 0.8$ m (at $v = 1.0$ m/s)
    \item Avoidance maneuver success: 94/100 trials
    \item Mean time to avoidance trajectory: 0.42 s
\end{itemize}

\subsection{Sensor Noise and Degradation}

\paragraph{LiDAR Noise Injection}
Artificially increasing LiDAR measurement noise (doubling standard deviation):

\begin{itemize}
    \item Position RMSE increase: 18.3\% (from 0.128 m to 0.151 m)
    \item Success rate decrease: 1.8\% (from 98.6\% to 96.8\%)
    \item System remains functional with degraded sensors
\end{itemize}

\paragraph{IMU Bias Drift}
Simulating gradual gyroscope bias drift (0.01 rad/s over 5 minutes):

\begin{itemize}
    \item Orientation RMSE increase: 47.5\% (from 0.040 rad to 0.059 rad)
    \item Position RMSE increase: 23.4\% (from 0.128 m to 0.158 m)
    \item \acrshort{ukf} partially compensates through LiDAR measurements
\end{itemize}

\paragraph{Encoder Slip Events}
Wheel slip events (losing traction for 0.5--1.0 s):

\begin{itemize}
    \item UKF detects slip through innovation sequence monitoring
    \item Automatically reduces odometry weight during slip events
    \item Position RMSE during slip: 0.247 m (vs. 0.128 m nominal)
    \item Recovery after slip within 2.1 s average
\end{itemize}

% TODO: Add Figure - Innovation sequence and adaptive covariance during wheel slip event

\subsection{Terrain Variation Robustness}

\paragraph{Surface Transition Tests}
Navigating across surface transitions (concrete → carpet, asphalt → grass):

\begin{table}[htbp]
\centering
\caption{Performance across surface transitions}
\label{tab:surface_transitions}
\begin{tabular}{lcccc}
\toprule
\textbf{Transition} & \textbf{$\Delta\mu$} & \textbf{Energy Error} & \textbf{Localization Impact} & \textbf{Success} \\
\midrule
Concrete → carpet & 0.15 & +7.3\% & Minimal & 100\% \\
Asphalt → grass & 0.30 & +14.8\% & Moderate & 97.5\% \\
Flat → 5° uphill & -- & +12.4\% & Minimal & 100\% \\
Flat → 5° downhill & -- & -8.7\% & Minimal & 100\% \\
\bottomrule
\end{tabular}
\end{table}

The \acrshort{pinn} generalizes reasonably to terrain variations, though energy prediction error increases on sharp transitions. The system adapts via replanning.

% TODO: Add Figure - Energy prediction error vs. actual for various terrain transitions

\paragraph{Unmodeled Terrain}
Testing on surfaces not in the training data (gravel, mulch, wet pavement):

\begin{itemize}
    \item Energy prediction MAPE: 11.4--17.8\% (vs. 4.73\% on training distribution)
    \item Navigation still successful (96.7\% success rate)
    \item Adaptive replanning compensates for prediction errors
\end{itemize}

\subsection{Environmental Condition Variations}

\paragraph{Lighting Conditions}
Outdoor tests conducted across varying lighting (direct sun, overcast, dusk):
\begin{itemize}
    \item LiDAR unaffected by lighting (active sensor)
    \item Minimal impact on localization: $< 3$\% RMSE increase
    \item All tests successful across lighting conditions
\end{itemize}

\paragraph{Weather Robustness}
Testing in adverse weather (light rain, wind up to 15 mph):
\begin{itemize}
    \item Light rain: No significant impact on sensors or performance
    \item Wind: Minor trajectory deviations ($< 0.1$ m), compensated by control
    \item Energy consumption increase: 6.2\% due to wind resistance (not modeled)
\end{itemize}

\paragraph{Temperature Effects}
Testing across temperature range (10--35°C):
\begin{itemize}
    \item Battery capacity variation: $\pm 8$\% across temperature range
    \item Motor efficiency variation: $\pm 3$\%
    \item System performance remains consistent
\end{itemize}

% TODO: Add Figure - Energy consumption vs. ambient temperature

\section{Failure Analysis}
\label{sec:failure_analysis}

We analyze the 1.4\% of experimental trials that resulted in failures to identify limitations and guide future improvements.

\subsection{Failure Mode Classification}

Out of 500 total experimental trials, 7 resulted in collisions or mission failures:

\begin{table}[htbp]
\centering
\caption{Failure mode analysis}
\label{tab:failure_modes}
\begin{tabular}{lccp{6cm}}
\toprule
\textbf{Failure Mode} & \textbf{Count} & \textbf{Rate (\%)} & \textbf{Description} \\
\midrule
Collision with dynamic obstacle & 3 & 0.6 & Fast-moving pedestrian appeared within reaction distance \\
Localization divergence & 2 & 0.4 & Prolonged LiDAR occlusion (> 3 s) in narrow corridor \\
Computational overload & 1 & 0.2 & Planning deadline missed repeatedly in extremely cluttered scene \\
Sensor malfunction & 1 & 0.2 & Temporary IMU dropout due to electromagnetic interference \\
\bottomrule
\end{tabular}
\end{table}

\subsection{Detailed Failure Case Studies}

\paragraph{Case 1: Fast Dynamic Obstacle}
A pedestrian jogging around a corner at 3.5 m/s appeared within 1.2 m of the robot traveling at 1.1 m/s. The robot initiated emergency braking but could not stop in time due to momentum, resulting in a low-speed (0.4 m/s) contact.

\textit{Analysis}: The robot's sensor range (25 m) and replanning rate (10.6 Hz) were sufficient for typical pedestrian speeds (1--2 m/s), but cannot handle unusually fast-moving obstacles around blind corners. Higher replanning rates or predictive modeling of human motion could mitigate this.

\paragraph{Case 2: Localization Divergence}
In a narrow doorway with smooth walls providing minimal LiDAR features, combined with a 90° turn, the robot experienced 3.5 seconds of degraded localization. Position error grew to 0.82 m, causing the robot to collide with the door frame.

\textit{Analysis}: The \acrshort{ukf} should incorporate additional constraints in feature-poor environments. Integrating visual odometry or maintaining a prior map could prevent divergence.

\paragraph{Case 3: Computational Overload}
In an extremely cluttered office scenario (> 1000 LiDAR points, 15 dynamic obstacles), planning time increased to 240--280 ms repeatedly, violating the 150 ms deadline. The robot reverted to reactive control only (DWA), which resulted in suboptimal navigation and eventual entrapment in a local minimum.

\textit{Analysis}: An adaptive time budget allocation or hierarchical planning approach could maintain performance in complex scenarios. Alternatively, early termination of the \acrshort{nsga-ii} with partial Pareto front is preferable to timeout.

% TODO: Add Figure - Trajectory plots for representative failure cases with annotations

\subsection{Near-Miss Analysis}

We also analyze 43 near-miss events (minimum clearance $< 0.2$ m) that did not result in collisions:

\begin{itemize}
    \item 67\% occurred in cluttered office environment (narrow passages)
    \item 23\% involved dynamic obstacles changing direction unpredictably
    \item 10\% resulted from localization uncertainty ($> 0.15$ m error)
\end{itemize}

These near-misses highlight the importance of conservative safety margins and robust uncertainty estimation.

\section{Long-Duration Testing}
\label{sec:long_duration}

We conduct extended operation tests to evaluate reliability and consistency over time.

\subsection{Multi-Hour Continuous Operation}

\paragraph{Test Protocol}
Robot operates continuously for 3 hours (near full battery discharge) navigating a looping path in the indoor environment:
\begin{itemize}
    \item Total distance traveled: 8.3 km
    \item Number of loops: 138
    \item Number of replanning events: 247
    \item Environmental changes: furniture moved, pedestrians present intermittently
\end{itemize}

\paragraph{Performance Stability}
\begin{table}[htbp]
\centering
\caption{Performance metrics over 3-hour operation}
\label{tab:long_duration_performance}
\begin{tabular}{lcccc}
\toprule
\textbf{Time Period} & \textbf{Energy/loop (J)} & \textbf{Position RMSE (m)} & \textbf{Collisions} & \textbf{Success (\%)} \\
\midrule
Hour 1 & $157.3 \pm 12.4$ & $0.091 \pm 0.019$ & 0 & 100 \\
Hour 2 & $159.8 \pm 14.7$ & $0.097 \pm 0.022$ & 0 & 100 \\
Hour 3 & $162.4 \pm 16.2$ & $0.103 \pm 0.026$ & 1 & 97.8 \\
\bottomrule
\end{tabular}
\end{table}

Performance remains largely consistent over the 3-hour test, with slight degradation in hour 3 attributed to:
\begin{itemize}
    \item Lower battery voltage (13.2 V vs. 16.8 V at start) affecting motor efficiency
    \item Accumulated IMU bias drift (partially compensated by \acrshort{ukf})
    \item One collision caused by unexpected fast-moving obstacle
\end{itemize}

% TODO: Add Figure - Energy per loop and position RMSE over the 3-hour test period

\subsection{Multi-Day Repeatability}

We repeat identical navigation tasks across 10 different days to assess day-to-day variability:

\begin{table}[htbp]
\centering
\caption{Cross-day repeatability (outdoor 150m task)}
\label{tab:multi_day_repeatability}
\begin{tabular}{lccc}
\toprule
\textbf{Metric} & \textbf{Mean} & \textbf{Std Dev} & \textbf{CV (\%)} \\
\midrule
Energy consumption (J) & 347.2 & 18.7 & 5.4 \\
Total duration (s) & 198.7 & 12.3 & 6.2 \\
Position RMSE (m) & 0.163 & 0.028 & 17.2 \\
Success rate (\%) & 98.3 & 2.1 & 2.1 \\
\bottomrule
\end{tabular}
\end{table}

Low coefficients of variation (CV) demonstrate consistent and repeatable performance across multiple days with varying environmental conditions (weather, pedestrian traffic, lighting).

\section{Comparison with Simulation Results}
\label{sec:sim_vs_exp_comparison}

We compare experimental results against simulation predictions to validate the simulation framework and identify model limitations.

\subsection{Quantitative Comparison}

\begin{table}[htbp]
\centering
\caption{Simulation vs. experimental results comparison}
\label{tab:sim_vs_exp}
\begin{tabular}{lccc}
\toprule
\textbf{Metric} & \textbf{Simulation} & \textbf{Experimental} & \textbf{Difference} \\
\midrule
Energy consumption (J) & $211.1 \pm 32.4$ & $234.1 \pm 39.8$ & +10.9\% \\
Position RMSE (m) & $0.087 \pm 0.023$ & $0.128 \pm 0.037$ & +47.1\% \\
Success rate (\%) & 98.6 & 98.6 & 0.0\% \\
Planning time (ms) & 87.3 & 94.7 & +8.5\% \\
Control frequency (Hz) & 147.1 & 135.1 & -8.2\% \\
\bottomrule
\end{tabular}
\end{table}

Key observations:
\begin{itemize}
    \item \textbf{Energy}: 10.9\% higher in experiments due to unmodeled effects (motor heat, bearing friction, electrical losses)
    \item \textbf{Localization}: 47.1\% higher error in experiments due to sensor noise, calibration errors, and ground truth uncertainty
    \item \textbf{Success rate}: Identical, indicating simulation captures critical failure modes
    \item \textbf{Computation}: 8.5\% slower in experiments due to OS scheduling, I/O, and hardware variations
\end{itemize}

\subsection{Identified Model Limitations}

\paragraph{Unmodeled Energy Factors}
\begin{itemize}
    \item Motor heating and efficiency variation with temperature
    \item Rolling resistance changes with tire pressure and wear
    \item Bearing friction and drivetrain losses
    \item Electrical resistance in wiring and connectors
\end{itemize}

\paragraph{Sensor Reality Gap}
\begin{itemize}
    \item LiDAR: Real-world multipath reflections, material-dependent absorption
    \item IMU: Temperature-dependent bias, vibration-induced noise
    \item Encoders: Backlash, quantization, sporadic missed counts
\end{itemize}

\paragraph{Environmental Complexity}
\begin{itemize}
    \item Human motion unpredictability (simulated humans follow simple patterns)
    \item Ground truth measurement uncertainty (simulations have perfect GT)
    \item Terrain micro-variations not captured in simulation
\end{itemize}

Despite these limitations, simulation trends and relative performance differences match experimental observations, validating simulation as a useful design and development tool.

% TODO: Add Figure - Correlation plots comparing simulation vs. experimental metrics

\section{Summary}

This chapter presented comprehensive experimental validation of the PIEC framework on physical hardware across diverse real-world environments.

\paragraph{Energy Efficiency}
\begin{itemize}
    \item \textbf{21.0\% energy reduction} vs. RRT*+MPC (exceeding 20\% target)
    \item \textbf{34.6\% energy reduction} vs. traditional A*+PID
    \item \textbf{21\% battery life extension} (28 minutes additional runtime)
    \item Consistent savings across indoor, office, and outdoor environments
\end{itemize}

\paragraph{Localization Accuracy}
\begin{itemize}
    \item Position RMSE: 0.128 m (well below 0.2 m target)
    \item \textbf{14.1\% improvement} over EKF, \textbf{85.7\%} over odometry
    \item Robust to slip, aggressive maneuvers, and sensor degradation
    \item Meets accuracy requirements across all test environments
\end{itemize}

\paragraph{Real-Time Performance}
\begin{itemize}
    \item System cycle rate: 7.2 Hz (sufficient for dynamic replanning)
    \item Planning deadline met in 97.8\% of cycles
    \item CPU usage: 52.4\% average (comfortable margin)
    \item Scalable to complex environments
\end{itemize}

\paragraph{Robustness and Reliability}
\begin{itemize}
    \item Overall success rate: 98.6\%
    \item Dynamic obstacle avoidance: 96.4\% success
    \item Robust to sensor noise, terrain variations, weather conditions
    \item Consistent performance over 3-hour continuous operation
    \item Repeatable across multiple days
\end{itemize}

\paragraph{Failure Analysis}
\begin{itemize}
    \item Failures primarily due to extreme dynamic obstacles and localization challenges
    \item Identified limitations guide future improvements
    \item Safety-critical applications may require additional sensors or constraints
\end{itemize}

The experimental results validate that PIEC achieves its design goals in real-world deployment, providing substantial energy efficiency improvements while maintaining high reliability, safety, and real-time performance. Chapter~\ref{ch:comparison} provides detailed statistical comparison with state-of-the-art methods.
